% Proof Notes for ccm::pow
% Compile with: pdflatex POW_PROOF.tex

\documentclass[11pt]{article}

\usepackage{amsmath}
\usepackage{amsthm}
\renewcommand{\qedsymbol}{\ensuremath{\blacksquare}}
\usepackage{amssymb}
\usepackage{geometry}
\usepackage{hyperref}
\usepackage{microtype}
\usepackage{lmodern}
\usepackage[T1]{fontenc}
\usepackage{booktabs}
\usepackage{algorithm}
\usepackage{algpseudocode}
\usepackage{listings}
\usepackage{xcolor}
\usepackage{enumitem}
\usepackage{longtable}

\geometry{margin=1.1in}
\hypersetup{colorlinks=true, linkcolor=blue, citecolor=blue, urlcolor=blue}

\lstset{
  basicstyle=\small\ttfamily,
  breaklines=true,
  columns=fullflexible,
  keepspaces=true,
  frame=single,
  backgroundcolor=\color{gray!6},
  rulecolor=\color{gray!40},
  framesep=4pt,
}

\theoremstyle{plain}
\newtheorem{theorem}{Theorem}
\newtheorem{lemma}{Lemma}
\newtheorem{corollary}{Corollary}
\newtheorem{definition}{Definition}

\theoremstyle{remark}
\newtheorem{remark}{Remark}

\newcommand{\ulp}{\operatorname{ulp}}
\newcommand{\Pow}{\texttt{ccm::pow}}
\newcommand{\Powf}{\texttt{ccm::powf}}
\newcommand{\Powl}{\texttt{ccm::powl}}
\newcommand{\Gdbl}{145805145804091023 \times 2^{-110}}
\newcommand{\Gflt}{577595370471038973 \times 2^{-111}}

\title{\textbf{Proof Notes for \texttt{ccm::pow}}}
\author{Ian Pike}
\date{}

\begin{document}
\maketitle
\thispagestyle{empty}

\begin{abstract}
These notes record proved statements for \Pow, \Powf, and \Powl.  For binary64
the input space has cardinality $2^{128}$, so the argument splits into separate
obligations on their natural domains, in the same style as CRlibm and
CORE-MATH~\cite{powr,coremath}: exceptional conformance, range-reduction
identities in exact rational arithmetic, approximation bounds with Sollya, and
rounded evaluation circuits with Gappa.  Where an end-to-end theorem needs a
transfer lemma that is not yet shown, the gap is stated explicitly.  This note
does not claim a complete binary64 $4$-ULP theorem on the whole non-exceptional
path.

Section~\ref{sec:powl80} covers the 80-bit \Powl{} kernel: Sollya sup-norm
bounds on the 64-bit significand, a double-double bound on the integer path,
and an explicit boundary on the reconstruction-circuit certificate.
\end{abstract}

\tableofcontents
\newpage

\section*{Notation}
\addcontentsline{toc}{section}{Notation}

For a floating-point format $F$, let $\mathbb{F}_F$ denote its representable
values.  In particular, $\mathbb{F}_{64}$ denotes IEEE~754 binary64 and
$\mathbb{F}_{32}$ denotes IEEE~754 binary32.

If $v$ is a nonzero finite number representable in $F$, then $\ulp_F(v)$
denotes the spacing of format $F$ at $v$.  In the normal range this is the
usual
\[
  \ulp_F(v) = 2^{\lfloor \log_2 |v| \rfloor - p + 1},
\]
where $p$ is the precision in bits.  In the subnormal range $\ulp_F(v)$ is
the constant minimum spacing.

Gappa~\cite{gappa} writes exact rationals in the form $N\mathtt{b}E$, meaning
$N \times 2^E$.

\section{Method and Scope}
\label{sec:scope}

\subsection{Objects under proof}

This note concerns the production sources listed below:
\begin{enumerate}[label=(O\arabic*)]
  \item \texttt{pow\_impl(double,double)} in \texttt{pow\_impl.hpp};
  \item \texttt{powf\_impl(float,float)} in \texttt{powf\_impl.hpp};
  \item the associated tables in \texttt{common\_math\_constants.hpp};
  \item the \Powl{} dispatch \texttt{powl\_impl} / \texttt{powl\_bounded\_integer}
    in \texttt{powl\_gen.hpp};
  \item the 80-bit kernel \texttt{powl\_ld80\_general\_finite} in
    \texttt{powl\_ld80\_kernel.hpp};
  \item the 80-bit tables and minimax coefficients in
    \texttt{powl\_ld80\_tables.hpp}.
\end{enumerate}
The supporting proof scripts read their constants directly from those files;
they do not rely on hand-copied duplicates.

The note separates exceptional semantics, range-reduction identities,
approximation on the reals, and rounded circuit evaluation; see~\cite{powr,coremath}.

\subsection{Standing assumptions}

Every non-exceptional theorem in this note is proved under the following
standing assumptions:
\begin{enumerate}[label=(A\arabic*)]
  \item IEEE~754 binary32 and binary64 semantics~\cite{ieee754};
  \item round-to-nearest, ties-to-even for the modeled operations;
  \item no flush-to-zero or denormals-are-zero mode;
  \item no fast-math reassociation or other algebraic transformation that
    changes the operation graph being certified;
  \item when a Gappa theorem is cited, the executed circuit is the non-FMA
    circuit encoded in the corresponding \texttt{.g} file;
  \item the \Powl{} discussion concerns exactly two platform shapes, and no
    others: the \emph{double-aliased} shape, with
    \texttt{numeric\_limits<long double>::digits=53}, on which \Powl{} is the
    binary64 kernel composed with the identity cast; and the
    \emph{x87 extended} shape, with \texttt{digits=64} and the macro
    \texttt{CCM\_TYPES\_LONG\_DOUBLE\_IS\_FLOAT80} defined, on which \Powl{} uses
    its own kernel and tables.  The IEEE binary128 and unknown shapes are out of
    scope here.
\end{enumerate}

\begin{remark}
Assumption~(A5) is substantive.  The present certificates are for explicit
non-FMA circuits.  They should not be read as blanket certificates for every
possible FMA-contracted build.
\end{remark}

\subsection{Reference value and ULP policy}

\begin{definition}[Reference value]
\label{def:reference}
For finite non-exceptional inputs $(x,y)$, the reference value is the real
number $x^y$, interpreted in the usual way:
\begin{enumerate}[label=(R\arabic*)]
  \item if $x>0$, then $x^y = 2^{y\log_2(x)}$;
  \item if $x<0$ and $y \in \mathbb{Z}$, then
    $x^y = \operatorname{sgn}(x)^y |x|^y$;
  \item exceptional cases are treated separately by the C/IEEE rules and do
    not belong to this definition.
\end{enumerate}
\end{definition}

\begin{definition}[4-ULP policy]
\label{def:ulp-policy}
Let $\widehat{r}$ be an output in format $F$, and let $r$ be the reference
value from Definition~\ref{def:reference}.
\begin{enumerate}[label=(U\arabic*)]
  \item If $r$ is finite and nonzero, the ULP error is
    $|\widehat{r} - r|/\ulp_F(r)$.
  \item If $r=0$, the requirement is the correct signed zero.
  \item If the C/IEEE rules require infinity, NaN, overflow, underflow, or a
    specific floating-point exception, then that mandated behavior replaces
    the finite ULP metric.
\end{enumerate}
\end{definition}

\subsection{What is proved here}

\begin{theorem}[Formal scope of this note]
\label{thm:scope}
Under the standing assumptions above, the present note establishes the
following statements, and no others.
\begin{enumerate}[label=(S\arabic*)]
  \item The exceptional branches of the binary64 kernel satisfy the stated
    C17/Annex~F cases.
  \item The first-stage and second-stage range reductions satisfy their stated
    interval bounds for every source-table entry and every representable input
    in the corresponding bucket.
  \item The source polynomials satisfy rigorous sup-norm bounds on their
    declared domains.
  \item The non-FMA floating-point evaluation circuits encoded in the Gappa
    files satisfy the displayed rounding-error bounds.
  \item Conditional on the residual variable $lo_6$ lying in
    $[-\tfrac12,\tfrac12]$, the exp stage of the binary64 kernel contributes
    at most $(G+S)2^{53} \le 1.025$ binary64 ULP, with $G=\Gdbl$ and $S$ the
    Sollya bound of Section~\ref{sec:poly}.
  \item For the binary32 primary path, the analogous exp-stage contribution is
    negligible compared with the final cast to float; the logical correctness
    of the Ziv acceptance test is proved separately in
    Theorem~\ref{thm:ziv}.
  \item On the x87 extended platform, the integer-exponent path of \Powl{}
    returns a value whose distance from $b^{n}$ is governed by the
    double-double error analysis of Theorem~\ref{thm:powl-int}.
  \item On the same platform, the two minimax polynomials of the general-finite
    \Powl{} kernel satisfy rigorous sup-norm bounds on the 64-bit significand
    (Lemmas~\ref{lem:supnorm-log-ld}--\ref{lem:supnorm-exp-ld}), and the
    kernel's range reduction inherits the interval bounds of
    Section~\ref{sec:reduction} verbatim.
  \item The two double-double polynomials of the accurate binary64
    reconstruction path satisfy rigorous sup-norm bounds on their declared
    intervals (Lemmas~\ref{lem:supnorm-log-dd}--\ref{lem:supnorm-exp-dd}), and
    the FMA error-free transformation used by the integer and half-integer
    algebraic path is exact in every IEEE rounding mode.
\end{enumerate}
No theorem stronger than \ref{thm:scope} is claimed in this document.  In
particular, no Gappa certificate is yet asserted for the 80-bit reconstruction
circuit or for the accurate binary64 double-double reconstruction circuit;
Section~\ref{sec:powl80} and Remark~\ref{rem:pow-accurate-scope} mark those
boundaries precisely.
\end{theorem}

\begin{remark}
Two transfer steps needed for a genuine end-to-end binary64 accuracy theorem
are still outside the present certificates.  First, the note does not yet
propagate the binary64 log-phase error all the way through the
\texttt{exact\_add}/\texttt{exact\_mult} reconstruction, the multiplication by
$y$, the scaling branch, and the final reconstruction of $2^{hm/64}$.  Second,
the note does not yet prove unconditionally that every non-exceptional binary64
execution reaches the exp polynomial with $|lo_6| \le \tfrac12$.  For that
reason Theorem~\ref{thm:exp-stage} is stated conditionally.
\end{remark}

\section{The Algorithm}

Both kernels are organized around the standard identity
\[
  x^y = 2^{\,y\log_2(x)}.
\]
The work is divided into an exceptional-input dispatch, a log phase, and an
exp phase.

The log phase writes the positive base in the form $x = 2^{e_x}m_x$ with
$m_x \in [1,2)$, chooses a range-reduction constant $r_i$ from the top seven
mantissa bits, and forms
\[
  dx = r_i m_x - 1.
\]
The tables are chosen so that $|dx| < 2^{-7}$.  Thus
\[
  \log_2(x) = e_x - \log_2(r_i) + \log_2(1+dx),
\]
and the nonlinear part of the problem has been confined to a very small
interval.

The exp phase forms an approximation to $64y\log_2(x)$, rounds its principal
part to an integer $hm$, reconstructs $2^{hm/64}$ from a table, and
approximates $2^{lo_6/64}$ by a short polynomial in the residual variable
$lo_6$.

\begin{algorithm}
\caption{Binary64 $\mathrm{pow}(x,y)$}
\label{alg:pow64}
\begin{algorithmic}[1]
\Require $x,y \in \mathbb{F}_{64}$
\Ensure The C/IEEE exceptional result, or a finite approximation to $x^y$

\Statex \textbf{Exceptional dispatch}
\If{$y=\pm 0$} \Return $1.0$ \EndIf
\If{$x=1$} \Return $1.0$ \EndIf
\If{$\mathrm{isnan}(x)$ or $\mathrm{isnan}(y)$} \Return qNaN \EndIf
\If{$x=\pm 0$} \Return signed $\pm\infty$ or signed $\pm0$ as required \EndIf
\If{$|x|=\infty$} \Return signed $\pm\infty$ or signed $\pm0$ as required \EndIf
\If{$|y|=\infty$} \Return $0$, $1$, or $\infty$ according to $|x| \lessgtr 1$ \EndIf
\If{$x<0$ and $y \notin \mathbb{Z}$} \Return qNaN and raise \texttt{FE\_INVALID} \EndIf
\State Record the sign determined by odd-integer parity; replace $x$ by $|x|$

\Statex \textbf{Algebraic exact routing}
\If{$y \in \mathbb{Z}$ and $|y| \le 1024$} \Return the double-double square-and-multiply value of \S\ref{sec:pow-exact} \EndIf
\If{$2y \in \mathbb{Z}$ and $|2y| \le 2048$} \Return $\sqrt{x}$ (for $y=\tfrac12$) or the double-double square-and-multiply of $\sqrt{x}$ raised to the odd integer $2y$ \EndIf

\Statex \textbf{Log phase}
\State Normalize subnormals if necessary
\State Write $x = 2^{e_x}m_x$ with $m_x \in [1,2)$
\State $i \gets \lfloor 128(m_x - 1) \rfloor$
\State $dx \gets \mathrm{RD}[i]m_x - 1$ or the algebraically equivalent compensated form
\State Reconstruct an approximation to $\log_2(x)$ from $e_x$, $\mathrm{LOG2\_R\_TD}[i]$, and a polynomial in $dx$

\Statex \textbf{Exp phase}
\State Form an approximation to $64y\log_2(x)$ as $T_{\mathrm{hi}} + T_{\mathrm{lo}}$
\State $hm \gets \mathrm{round}(T_{\mathrm{hi}})$
\State $lo_6 \gets T_{\mathrm{hi}} - hm + T_{\mathrm{lo}}$
\State Reconstruct $2^{hm/64}$ from the table \texttt{EXP2\_MID1}
\State Evaluate a polynomial approximation to $2^{lo_6/64}$
\If{the result is not clamped to overflow or underflow}
  \State Carry the log and exp phases in double-double and \Return one binary64 rounding of $r_{\mathrm{hi}} + r_{\mathrm{lo}}$ (\S\ref{sec:pow-accurate})
\Else{}
  \State \Return the saturating single-double product, with the recorded sign and scaling factor
\EndIf
\end{algorithmic}
\end{algorithm}

\begin{algorithm}
\caption{Binary32 $\mathrm{powf}(x,y)$}
\label{alg:powf}
\begin{algorithmic}[1]
\Require $x,y \in \mathbb{F}_{32}$
\Ensure The C/IEEE exceptional result, or a float approximation to $x^y$

\State Exceptional dispatch and short fast paths
\State Perform the primary log/exp computation in binary64 working precision
\State Apply the Ziv test
\If{the $\pm64$-ULP-double neighborhood rounds to one float}
  \Return that float
\EndIf
\State Otherwise enter the double-double fallback and round its result to float
\end{algorithmic}
\end{algorithm}

\section{Exceptional Inputs}
\label{sec:exceptions}

\begin{theorem}[Exceptional branch correctness]
\label{thm:exceptional}
For binary64 inputs, the exceptional branches of
\texttt{pow\_impl(double,double)} cover the cases required by
C17 \S7.12.7.4 and Annex~F, \S F.10.4.4~\cite{c17}, subject to the standing
assumptions of Section~\ref{sec:scope}.
\end{theorem}

\begin{proof}
The guards are evaluated in a fixed order, and that order matters.  The two
identities
\[
  \mathrm{pow}(\mathrm{NaN},0)=1,
  \qquad
  \mathrm{pow}(1,\mathrm{NaN})=1
\]
must be recognized before generic NaN propagation; the implementation does so.
After those two special cases, the remaining branches are disjoint and cover
the classical exceptional cases.

\begin{center}
\begin{tabular}{@{}p{0.28\linewidth}p{0.29\linewidth}p{0.30\linewidth}@{}}
\toprule
Guard family & Result & Reason \\
\midrule
\texttt{exp == 0}
  & $1.0$
  & $\mathrm{pow}(x,\pm0)=1$ for all $x$. \\[4pt]
\texttt{base == 1}
  & $1.0$
  & $\mathrm{pow}(1,y)=1$ for all $y$. \\[4pt]
\texttt{isnan(base)} or \texttt{isnan(exp)}
  & qNaN
  & NaN propagation after the two fixed points above. \\[4pt]
\texttt{base == +-0}, \texttt{exp < 0}
  & signed $\pm\infty$
  & Pole case; \texttt{FE\_DIVBYZERO} is raised. \\[4pt]
\texttt{base == +-0}, \texttt{exp > 0}
  & signed $\pm0$
  & Sign determined by odd-integer parity. \\[4pt]
\texttt{base == -inf}, non-integer \texttt{exp}
  & qNaN
  & Domain error; \texttt{FE\_INVALID} is raised. \\[4pt]
\texttt{base == +-inf}, integer \texttt{exp}
  & signed $\pm\infty$ or signed $\pm0$
  & Determined by sign, parity, and the sign of the exponent. \\[4pt]
\texttt{exp == +-inf}
  & $0$, $1$, or $\infty$
  & Determined by comparing $|x|$ with $1$. \\[4pt]
finite \texttt{base < 0}, non-integer \texttt{exp}
  & qNaN
  & Domain error; \texttt{FE\_INVALID} is raised. \\
\bottomrule
\end{tabular}
\end{center}

The auxiliary predicates \texttt{is\_integer} and \texttt{is\_odd\_integer}
operate directly on the binary64 exponent and significand structure.  For
magnitudes at least $2^{53}$ every finite binary64 value is already an integer,
so the parity question is reduced to the least significant stored unit bit.
Thus the large-exponent cases do not require a separate argument.
\end{proof}

\section{Range Reduction}
\label{sec:reduction}

\begin{lemma}[First-stage range reduction]
\label{lem:stage1}
Let $x$ be a finite positive binary64 number with mantissa $m_x \in [1,2)$,
and let
\[
  i = \lfloor 128(m_x-1) \rfloor.
\]
Then $i \in \{0,\dots,127\}$, and the following two facts hold.
\begin{enumerate}[label=(\roman*)]
  \item The source-table identity
    \[
      \mathrm{RD}[i]\!\left(1 + \frac{i}{128}\right) - 1 = \mathrm{CD}[i]
    \]
    holds exactly.
  \item For every representable $m_x$ in bucket $i$,
    \[
      dx = \mathrm{RD}[i]m_x - 1 \in [-2^{-8},\,2^{-7}).
    \]
\end{enumerate}
\end{lemma}

\begin{proof}
The proof script \texttt{pow\_double\_proof.py} parses the arrays
\texttt{RD} and \texttt{CD} directly from the production header and checks
every bucket in exact rational arithmetic using Python's \texttt{Fraction}.

Part~(i) is a direct equality test.  For part~(ii), the script evaluates the
two bucket endpoints
\[
  1 + \frac{i}{128},
  \qquad
  1 + \frac{i+1}{128} - 2^{-52},
\]
passes them through the affine map $t \mapsto \mathrm{RD}[i]t - 1$, and notes
that a positive-slope affine map carries the whole interval between the images
of its endpoints.  Since the 128 buckets cover the entire mantissa range
$[1,2)$, the result follows for every finite positive binary64 input.
\end{proof}

\begin{remark}
The analogous binary32 statement is verified by
\texttt{pow\_float\_proof.py}; in addition, that script checks the exact
identity \texttt{R[k] == RD[k]} after promotion from float to double.
\end{remark}

\begin{lemma}[Second-stage range reduction]
\label{lem:stage2}
In the \Powf{} double-double fallback, once the first stage has produced
$dx \in [-2^{-8},2^{-7})$, the second reduction satisfies
\[
  dx_2 = (1+dx)\mathrm{R2}[i] - 1
  \in [-\mathtt{0x1.3ffcp-15},\,\mathtt{0x1.3e3dp-15}].
\]
\end{lemma}

\begin{proof}
The proof is again exhaustive over the source table.  For each entry of
\texttt{R2}, the script intersects the admissible first-stage interval with
the bucket corresponding to that entry, and then maps the resulting interval
through the affine function $u \mapsto (1+u)\mathrm{R2}[i]-1$.  Everything is
done in exact rational arithmetic, so the displayed bound is exact as stated.
\end{proof}

\section{Polynomial Approximation}
\label{sec:poly}

Once the range reduction has been established, the approximation problem is
small enough to be treated by short source polynomials on fixed intervals.

\begin{lemma}[Log polynomial bound]
\label{lem:supnorm-log}
Let $P_{\!\log}$ be the binary64 polynomial with coefficients
\texttt{POW\_LOG2\_COEFFS}.  Then
\[
  \sup_{dx \in [-2^{-8},\,2^{-7}]}
  \left|
    P_{\!\log}(dx) - \frac{\log_2(1+dx)}{dx}
  \right|
  \leq S_{\log},
\]
where $S_{\log}$ is the hexadecimal upper bound emitted by Sollya.
\end{lemma}

\begin{proof}
The proof script feeds the source polynomial directly to Sollya and evaluates
\[
  \texttt{dirtyinfnorm(log2(1 + x)/x - Plog, [-2\string^-8; 2\string^-7])}.
\]
Sollya computes a rigorous interval-arithmetic upper bound for the sup norm,
not a sampled estimate.
\end{proof}

\begin{lemma}[Exp polynomial bound]
\label{lem:supnorm-exp}
Let $P_{\!\exp}$ be the binary64 polynomial with coefficients
\texttt{POW\_EXP2\_COEFFS}.  Then
\[
  \sup_{z \in [-1/2,\,1/2]}
  |P_{\!\exp}(z) - 2^{z/64}|
  \leq S_{\exp},
\]
where $S_{\exp}$ is the hexadecimal upper bound emitted by Sollya.
\end{lemma}

\begin{proof}
This is the same argument, applied to
\[
  \texttt{dirtyinfnorm(2\string^(x/64) - Pexp, [-2\string^-1; 2\string^-1])}.
\]
\end{proof}

\section{Floating-Point Evaluation Certificates}
\label{sec:gappa}

The preceding section concerns the exact real polynomials.  The implementation
evaluates them in floating-point arithmetic; the certificates below are Gappa
bounds on the encoded non-FMA circuits.

\subsection{Mapping from source code to proof model}

The Gappa files do not certify an abstract polynomial; they certify explicit
rounded circuits.  The relevant correspondences are listed below.  In every
row, \texttt{rnd()} denotes one round-to-nearest-even binary64 operation.

\begin{longtable}{@{}p{0.29\linewidth}p{0.41\linewidth}p{0.12\linewidth}p{0.11\linewidth}@{}}
\toprule
Source step & Gappa model & Mode & File \\
\midrule
\endhead
\texttt{dx2 = dx * dx}
  & \texttt{x2 = rnd(x * x)}
  & RN & \texttt{pow\_double\_log\_eval.g} \\[4pt]
\texttt{c0 = muladd(dx, C2, C1)}
  & \texttt{m0 = rnd(x * C2);} \texttt{ c0 = rnd(m0 + C1)}
  & RN & same \\[4pt]
\texttt{c1 = muladd(dx, C4, C3)}
  & \texttt{m1 = rnd(x * C4);} \texttt{ c1 = rnd(m1 + C3)}
  & RN & same \\[4pt]
\texttt{c2 = muladd(dx, C6, C5)}
  & \texttt{m2 = rnd(x * C6);} \texttt{ c2 = rnd(m2 + C5)}
  & RN & same \\[4pt]
\texttt{p = polyeval(dx2, c0, c1, c2)}
  & \texttt{h0,h1,h2,h3,h4}
  & RN & same \\[4pt]
\texttt{log polynomial output}
  & \texttt{log\_eval = rnd(h4 + x * C0)}
  & RN & same \\[8pt]

\texttt{lo6\_sqr = lo6 * lo6}
  & \texttt{x2 = rnd(x * x)}
  & RN & \texttt{pow\_double\_exp\_eval.g} \\[4pt]
\texttt{d0 = muladd(lo6, C2, C1)}
  & \texttt{m0 = rnd(x * C2);} \texttt{ d0 = rnd(m0 + C1)}
  & RN & same \\[4pt]
\texttt{d1 = muladd(lo6, C4, C3)}
  & \texttt{m1 = rnd(x * C4);} \texttt{ d1 = rnd(m1 + C3)}
  & RN & same \\[4pt]
\texttt{pp = polyeval(lo6\_sqr, d0, d1, C5)}
  & \texttt{h0,h1,h2,h3,h4}
  & RN & same \\[4pt]
\texttt{exp polynomial output}
  & \texttt{exp\_eval = rnd(h4 + 0x1p0)}
  & RN & same \\[8pt]

\texttt{powf} primary log circuit
  & \texttt{m0/c0, m1/c1, m2/c2, h0,h1,h2}
  & RN & \texttt{pow\_float\_log\_eval.g} \\[4pt]
\texttt{powf} primary exp circuit
  & \texttt{m0/d0, m1/d1, m2/d2, h0,h1,h2}
  & RN & \texttt{pow\_float\_exp\_eval.g} \\
\bottomrule
\end{longtable}

\begin{remark}
This table is the precise boundary of the Gappa argument.  If a build changes
the circuit, the proof file must change with it.
\end{remark}

\subsection{Certified rounding-error bounds}

\begin{lemma}[Binary64 log evaluation certificate]
\label{lem:gappa-dbl-log}
For the non-FMA binary64 log circuit,
\[
  |\mathrm{log\_eval} - P| \leq 584853814232881131 \times 2^{-119}.
\]
\end{lemma}

\begin{lemma}[Binary64 exp evaluation certificate]
\label{lem:gappa-dbl-exp}
For the non-FMA binary64 exp circuit,
\[
  |\mathrm{exp\_eval} - \mathrm{exp\_stagee}| \leq \Gdbl.
\]
\end{lemma}

\begin{lemma}[Binary32 log evaluation certificate]
\label{lem:gappa-flt-log}
For the non-FMA binary64 working-precision log circuit used in the primary
path of \Powf,
\[
  |\mathrm{log\_eval} - P| \leq 577604123977524393 \times 2^{-111}.
\]
\end{lemma}

\begin{lemma}[Binary32 exp evaluation certificate]
\label{lem:gappa-flt-exp}
For the corresponding exp circuit,
\[
  |\mathrm{exp\_eval} - P| \leq \Gflt.
\]
\end{lemma}

\begin{proof}[Proof of Lemmas \ref{lem:gappa-dbl-log}--\ref{lem:gappa-flt-exp}]
Each \texttt{.g} file introduces the rounded circuit and its exact real
counterpart side by side, then asks Gappa to prove a symmetric interval bound
for the difference.  The proof scripts treat failure as fatal and can archive
the command line, source hashes, and tool output in a JSON report.  Since the
reported intervals are the very intervals used by Gappa, the lemmas are simply
restatements of those certificates.
\end{proof}

\section{Consequences for the Exp Stage}
\label{sec:exp}

The exp polynomial is certified on the interval $[-\tfrac12,\tfrac12]$, and
the theorem below says exactly what follows when the implementation reaches the
polynomial inside that interval.

\begin{theorem}[Conditional binary64 exp-stage bound]
\label{thm:exp-stage}
Assume a finite non-exceptional binary64 input pair, and assume further that
the residual variable supplied to the exp polynomial satisfies
$|lo_6| \le \tfrac12$.  Then the exp stage contributes at most
\[
  (G+S)\,2^{53}
\]
binary64 ULP, where $G=\Gdbl$ is the Gappa bound of
Lemma~\ref{lem:gappa-dbl-exp} and $S$ is the Sollya bound of
Lemma~\ref{lem:supnorm-exp}.  Numerically this is at most $1.025$ binary64
ULP.
\end{theorem}

\begin{proof}
Lemma~\ref{lem:gappa-dbl-exp} gives
\[
  |\mathrm{exp\_eval} - P_{\!\exp}(lo_6)| \le G,
\]
while Lemma~\ref{lem:supnorm-exp} gives
\[
  |P_{\!\exp}(lo_6) - 2^{lo_6/64}| \le S.
\]
The triangle inequality yields
\[
  |\mathrm{exp\_eval} - 2^{lo_6/64}| \le G+S.
\]
Multiplying by the reconstructed factor $2^{hm/64}$ produces an absolute
output error bounded by $2^{hm/64}(G+S)$.  When this is measured in ULP at the
true result, the worst case occurs when $hm$ is a multiple of $64$ and
$lo_6<0$, which introduces one extra factor of $2$.  Hence the contribution is
bounded by $(G+S)2^{53}$ binary64 ULP.
\end{proof}

\begin{corollary}[Binary32 exp-stage contribution]
\label{cor:float-exp}
For the primary path of \Powf, the analogous argument yields
\[
  (G_f + S_f)\,2^{24} \approx 3.76 \times 10^{-9}
  \quad\text{float ULP},
\]
which is negligible in comparison with the final cast to float.
\end{corollary}

\section{Accurate Binary64 Reconstruction}
\label{sec:pow-accurate}

For finite non-exceptional inputs that do not clamp to overflow or underflow,
the binary64 kernel reconstructs $x^y$ in double-double working precision and
returns one binary64 rounding of $r_{\mathrm{hi}} + r_{\mathrm{lo}}$.  The
single-double circuit of Sections~\ref{sec:poly}--\ref{sec:gappa} is retained
only on the clamped over/underflow tail, where the result saturates to a signed
infinity or zero.

The accurate path repeats the second-stage range reduction of
Lemma~\ref{lem:stage2} in double-double, evaluates a degree-5 double-double
minimax for $\log_2(1+\cdot)$ and a degree-9 double-double minimax for
$2^{\cdot/64}$, and collapses the final pair with one binary64 addition.  The
two minimax polynomials are the new numeric artifacts; their approximation
error is certified below.

\begin{lemma}[Accurate log2 double-double polynomial bound]
\label{lem:supnorm-log-dd}
Let $P_{\!\log}^{\mathrm{dd}}$ be the degree-5 polynomial whose coefficients are
the double-double array \texttt{COEFFS} in \texttt{pow\_impl.hpp}, each summed as
$\mathit{hi}+\mathit{lo}$.  Then
\[
  \sup_{x \in [-\mathtt{0x1.3ffcp-15},\,\mathtt{0x1.3e3dp-15}]}
  \left| P_{\!\log}^{\mathrm{dd}}(x) - \frac{\log_2(1+x)}{x} \right|
  \;\le\; S_{\log}^{\mathrm{dd}},
\]
where $S_{\log}^{\mathrm{dd}} \approx 2^{-96}$ is the rigorous sup-norm bound
emitted by Sollya.
\end{lemma}

\begin{lemma}[Accurate exp2 double-double polynomial bound]
\label{lem:supnorm-exp-dd}
Let $P_{\!\exp}^{\mathrm{dd}}$ be the degree-9 polynomial whose coefficients are
the double-double array \texttt{EXP2\_COEFFS} in \texttt{pow\_impl.hpp}, each
summed as $\mathit{hi}+\mathit{lo}$.  Then
\[
  \sup_{z \in [-1/2,\,1/2]}
  \left| P_{\!\exp}^{\mathrm{dd}}(z) - 2^{z/64} \right|
  \;\le\; S_{\exp}^{\mathrm{dd}},
\]
where $S_{\exp}^{\mathrm{dd}} \approx 2^{-106}$ is the rigorous sup-norm bound
emitted by Sollya.
\end{lemma}

\begin{proof}[Proof of Lemmas \ref{lem:supnorm-log-dd} and \ref{lem:supnorm-exp-dd}]
The script \texttt{pow\_double\_proof.py} parses the double-double arrays from
\texttt{pow\_impl.hpp}, reconstructs each coefficient as the exact sum of its two
limbs, and asks Sollya for \texttt{dirtyinfnorm} of each residual on its declared
interval.  As in Section~\ref{sec:poly} the returned magnitude is a rigorous
interval bound, not a sample.  The script fails unless the bounds stay below
$2^{-90}$ and $2^{-95}$ respectively, well under the $2^{-52}$ needed for a
single binary64 rounding to be faithful.
\end{proof}

\subsection{Exact integer and half-integer exponents}
\label{sec:pow-exact}

Integer-exponent results sit within roughly $2^{-106}$ of a representable double
and half-integer results land on exact midpoints, so the log/exp round-trip
cannot resolve their rounding from the active mode alone.  These exponents are
routed to an algebraic path before the kernel.

For integer $y$ with $|y| \le 1024$ the kernel computes $x^{|y|}$ by
double-double square-and-multiply.  Each product uses the FMA error-free
transformation
\[
  \mathit{hi} = \mathrm{fl}(ab),
  \qquad
  \mathit{lo} = \mathrm{fma}(a,b,-\mathit{hi}),
\]
which is exact in every IEEE rounding mode, and a \textsc{FastTwoSum}
renormalization.  For negative $y$ one Newton reciprocal of the resulting pair is
taken.  For half-integer $y$ with $|2y| \le 2048$ the identity
$x^{y} = (\sqrt{x})^{2y}$ is used: $y=\tfrac12$ returns the correctly-rounded
\texttt{sqrt\_gen}, and otherwise the kernel raises a double-double $\sqrt{x}$ to
the odd integer $2y$ by the same square-and-multiply. If the algebraic result is
not a finite normal double, the input is handed back to the kernel so that the
overflow, underflow, and subnormal semantics of the standard are preserved.

\begin{remark}
\label{rem:pow-accurate-scope}
The certificates of this section cover the approximation error of the two
double-double polynomials (Lemmas~\ref{lem:supnorm-log-dd}
and~\ref{lem:supnorm-exp-dd}) and the mode-independent exactness of the FMA
error-free transformation used by the algebraic path.  They do not constitute a
Gappa certificate for the full rounded double-double reconstruction circuit, and
no bit-exact correct-rounding theorem is asserted here.  End-to-end
correct-rounding in all four IEEE modes is validated empirically by the
CORE-MATH all-mode campaign over the generic runtime path; that measurement is
engineering evidence and is not part of the formal certificates.
\end{remark}

\section{The Ziv Test for \Powf}
\label{sec:ziv}

\begin{theorem}[Correctness of the Ziv acceptance predicate]
\label{thm:ziv}
Let $r$ be a binary64 number, and let
\[
  \rho_+ = \mathrm{fl}_{32}(r + 64\,\ulp_{64}(r)),
  \qquad
  \rho_- = \mathrm{fl}_{32}(r - 64\,\ulp_{64}(r)).
\]
If $\rho_+ = \rho_-$, then every real number in the interval
\[
  [\,r - 64\,\ulp_{64}(r),\; r + 64\,\ulp_{64}(r)\,]
\]
rounds to that same float under round-to-nearest, ties-to-even.
\end{theorem}

\begin{proof}
The rounding map $\mathrm{fl}_{32}$ is monotone on the real line.  Therefore,
if the two endpoints of a closed interval round to the same float, then every
point between them must round to that same float.  This is exactly the test
performed by the primary path of \Powf.
\end{proof}

\begin{corollary}
If the true real result of the primary path lies within
$\pm 64\,\ulp_{64}(r)$ of the computed binary64 value $r$, and if the Ziv
predicate accepts, then the returned float is correctly rounded.
\end{corollary}

\section{The 80-bit \texttt{powl} Kernel}
\label{sec:powl80}

When \texttt{long double} is the x87 extended type with a 64-bit significand,
\Powl{} uses a dedicated kernel.

\subsection{Dispatch}

Exceptional cases match Theorem~\ref{thm:exceptional}, except that a negative
infinite base follows glibc's parity convention rather than the textbook one.
For a finite positive base and finite exponent $y$, when
\texttt{try\_extract\_int64} recovers an integer with $|y| \le 2^{62}-1$, the
kernel uses \S\ref{sec:powl-int}; otherwise it uses \S\ref{sec:powl-gen}.

\subsection{Integer exponents by compensated squaring}
\label{sec:powl-int}

The integer path is square-and-multiply with double-double accumulators for the
running product and square.

\begin{algorithm}
\caption{80-bit integer power $b^{n}$}
\label{alg:powl-int}
\begin{algorithmic}[1]
\Require finite nonzero $b$; integer $n$ with $1 \le |n| \le 2^{62}-1$
\State $P \gets (1,0)$; \quad $F \gets (b,0)$; \quad $e \gets |n|$
\While{$e > 0$}
  \If{$e$ is odd} $P \gets P \otimes F$ \EndIf
  \State $e \gets \lfloor e/2 \rfloor$
  \If{$e > 0$} $F \gets F \otimes F$ \EndIf
\EndWhile
\If{$n < 0$} \Return one Newton reciprocal of the pair $P$
\Else{} \Return $\mathrm{round}_{80}(P_{\mathrm{hi}} + P_{\mathrm{lo}})$
\EndIf
\end{algorithmic}
\end{algorithm}

The operator $\otimes$ is a guarded double-double product.  It multiplies the
high limbs with an exact \textsc{TwoProduct}, folds in the cross terms, and
renormalizes with a \textsc{FastTwoSum}.  The one exception is that when either operand is near \texttt{LDBL\_MAX} it
forms the plain product and sets the residual to zero, as noted below.

\begin{theorem}[Compensated integer power]
\label{thm:powl-int}
Let $b$ be a finite nonzero long double and $n$ an integer with
$1 \le |n| \le 2^{62}-1$, and suppose every partial product $b^{2^{j}}$ and every
running partial result encountered by Algorithm~\ref{alg:powl-int} is a normal
long double of magnitude below the split-safety threshold $2^{16300}$.  Then the
pair $(P_{\mathrm{hi}}, P_{\mathrm{lo}})$ produced by the loop satisfies
\[
  \bigl| (P_{\mathrm{hi}} + P_{\mathrm{lo}}) - b^{|n|} \bigr|
  \;\le\; \eta\,\bigl| b^{|n|} \bigr|,
  \qquad \eta \le 2^{-110}.
\]
Consequently $\mathrm{round}_{80}(P_{\mathrm{hi}} + P_{\mathrm{lo}})$ differs from
$b^{|n|}$ by at most $\tfrac12 + 2^{-45}$ ULP.  For $n < 0$ the value returned is
one compensated reciprocal of the pair, which adds at most one further ULP.  In
every case the integer path meets the 4-ULP policy of
Definition~\ref{def:ulp-policy}.
\end{theorem}

\begin{proof}
The loop runs at most $2\lfloor \log_2 |n| \rfloor \le 124$ double-double
multiplications, one square for each bit of $|n|$ and one product for each bit
that is set.  Every one of them is the standard double-double multiplication: an
exact \textsc{TwoProduct} (the Veltkamp split \texttt{exact\_mult}, exact so long
as its operand stays below $2^{16383-33}$), then a normalizing
\textsc{FastTwoSum}.  Its relative error is at most $c\,u^{2}$, where
$u = 2^{-64}$ and $c$ is the small constant of the textbook
analysis~\cite{handbook}.  Treat each step as multiplication by $1 + \delta_k$
with $|\delta_k| \le c\,u^{2}$; composing $k \le 124$ of them gives
\[
  \prod_{k}(1 + \delta_k) - 1
  \;\le\; (1 + c\,u^{2})^{124} - 1
  \;\le\; 2^{-110},
\]
the $\eta$ that was claimed.  Now $\eta$ falls short of $2^{-64}$ by more than
forty binary orders, so $P_{\mathrm{lo}}$ carries the residual well below the
64th bit, and the only error of any size in the returned scalar is the final
rounding of $P_{\mathrm{hi}} + P_{\mathrm{lo}}$, which costs half an ULP.  When
$n$ is negative the kernel inverts the pair with one Newton step; starting from a
pair this good, a single step lands inside one ULP.  The hypotheses on normal
range and split safety are just the conditions under which the guard in
\texttt{powl\_dd::mul} keeps the Veltkamp constant from overflowing.  Outside
them the kernel stops compensating and returns a saturating infinity or zero,
the overflow or underflow result required by the standard.
\end{proof}

\begin{remark}
The threshold $2^{16300}$ is required because the Veltkamp split scales its
argument by $2^{33}$.  Near \texttt{LDBL\_MAX} that scaled value is already
infinite, so the following subtraction $\infty - \infty$ would yield NaN in the
residual.  The guard returns a clean infinity instead.  The hypothesis is the
range where \texttt{powl\_dd::mul} keeps the split exact.
\end{remark}

\subsection{The general-finite kernel}
\label{sec:powl-gen}

The general-finite path follows Algorithm~\ref{alg:pow64} with double-double
intermediates.  Range reduction matches the binary64 case: the same seven
mantissa bits choose a bucket, the same \texttt{RD} and \texttt{CD} entries form
$dx = \mathrm{RD}[i] m_x - 1$, and those entries are the binary64 constants
promoted to long double without rounding error.

\begin{corollary}[80-bit range reduction]
\label{cor:ld-reduction}
For every finite positive long double $x$ the reduced argument satisfies
$dx \in [-2^{-8}, 2^{-7})$.
\end{corollary}

\begin{proof}
Lemma~\ref{lem:stage1} goes through word for word.  The table and the bucket map
are identical, and a positive-slope affine map sends an interval to the interval
between the images of its endpoints.  The only thing that moves is the right end
of each bucket, from $1 + (i+1)/128 - 2^{-52}$ to $1 + (i+1)/128 - 2^{-63}$, and
the new endpoint sits inside the old one, hence inside the same image.
\end{proof}

The logarithm is built from $e_x$, the binary64 triple-double table
\texttt{LOG2\_R\_TD} (read back as long double, again exactly), and the log
polynomial of \S\ref{sec:powl-poly}.  The exponential stage forms
$64\,y\log_2 x$ as a pair $(T_{\mathrm{hi}}, T_{\mathrm{lo}})$, takes
$hm = \mathrm{round}(T_{\mathrm{hi}})$ and
$lo_6 = (T_{\mathrm{hi}} - hm) + T_{\mathrm{lo}}$, rebuilds $2^{hm/64}$ from the
\texttt{EXP2} table as a pair $(e_{\mathrm{hi}}, e_{\mathrm{lo}})$, and evaluates
the exp polynomial at $lo_6$.

Two reconstruction steps matter at 64-bit precision:
\begin{enumerate}[label=(\roman*)]
  \item \emph{Take the integer part from the high limb alone.}  With
    $hm = \mathrm{round}(T_{\mathrm{hi}})$ the difference $T_{\mathrm{hi}} - hm$
    is exact, by Sterbenz's lemma, so $lo_6 = (T_{\mathrm{hi}} - hm) +
    T_{\mathrm{lo}}$ keeps full precision.  Had one instead added
    $T_{\mathrm{hi}}$ and $T_{\mathrm{lo}}$ into a single long double first, the
    sum would round at the magnitude of $hm$ and throw away about $\log_2|hm|$
    low bits of $lo_6$.  At binary64 the loss hides beneath the result; at 64
    bits it does not.
  \item \emph{Scale the low limb as well.}  The pair
    $e_{\mathrm{hi}} + e_{\mathrm{lo}}$ holds $2^{hm/64}$ to nearly twice the
    significand width, so its product with
    $2^{lo_6/64} = 1 + lo_6\,P_{\!\exp}^{80}(lo_6)$ owes a term
    $e_{\mathrm{lo}}\,(2^{lo_6/64} - 1)$.  Omit it, as a binary64 kernel may with
    impunity, and about $2^{-60}$ of relative error escapes; that is nothing at
    binary64, and several long double ULP once the mid-table entry is nonzero and
    $lo_6$ approaches $\tfrac12$.
\end{enumerate}

When $|T_{\mathrm{hi}}|$ runs past $512 \cdot 64$ the kernel peels off multiples
of $512 \cdot 64$ before rounding and puts the matching power of two back at the
end.  This is the $2^{\pm 512}$ trick of the binary64 kernel, repeated as often
as the exponent demands, so that every finite power within range is reached.

\subsection{The 80-bit minimax certificates}
\label{sec:powl-poly}

The 80-bit kernel uses higher-degree minimax polynomials than binary64.  On the
intervals here the best errors at the binary64 degrees are about $2^{-56}$ for
the logarithm and $2^{-60}$ for the exponential, under one binary64 ULP but
several ULP at 64 bits.  The refit adds one term to each polynomial.

\begin{lemma}[80-bit log polynomial bound]
\label{lem:supnorm-log-ld}
Let $P_{\!\log}^{80}$ be the degree-7 polynomial with coefficients
\texttt{POW\_LOG2\_COEFFS} of \texttt{powl\_ld80\_tables.hpp}.  Then
\[
  \sup_{dx \in [-2^{-8},\,2^{-7}]}
  \left| P_{\!\log}^{80}(dx) - \frac{\log_2(1 + dx)}{dx} \right|
  \;\le\; S_{\log}^{80},
\]
where $S_{\log}^{80} \approx 2^{-65}$ is the rigorous sup-norm bound emitted by
Sollya.
\end{lemma}

\begin{lemma}[80-bit exp polynomial bound]
\label{lem:supnorm-exp-ld}
Let $P_{\!\exp}^{80}$ be the degree-6 polynomial with coefficients
\texttt{POW\_EXP2\_COEFFS} of \texttt{powl\_ld80\_tables.hpp}.  Then
\[
  \sup_{z \in [-1/2,\,1/2]}
  \left| P_{\!\exp}^{80}(z) - 2^{z/64} \right|
  \;\le\; S_{\exp}^{80},
\]
where $S_{\exp}^{80} \approx 2^{-71}$ is the rigorous sup-norm bound emitted by
Sollya.
\end{lemma}

\begin{proof}[Proof of Lemmas \ref{lem:supnorm-log-ld} and \ref{lem:supnorm-exp-ld}]
The script \texttt{validate\_powl\_ld80\_kernel.sollya} carries the production
coefficients and asks Sollya for \texttt{dirtyinfnorm} of each residual on its
interval.  As in Section~\ref{sec:poly}, what comes back is a rigorous interval
bound on the sup norm rather than a sample; the magnitudes shown are those
bounds, rounded up.
\end{proof}

\subsection{What is proved, and what is not}

Proved for the general-finite \Powl{} path: exceptional dispatch
(Theorem~\ref{thm:exceptional}), range reduction (Corollary~\ref{cor:ld-reduction}),
polynomial bounds (Lemmas~\ref{lem:supnorm-log-ld} and~\ref{lem:supnorm-exp-ld}),
and the integer path (Theorem~\ref{thm:powl-int}).  Not yet proved: a Gappa
certificate for the reconstruction circuit (\texttt{exact\_mult},
\texttt{exact\_add}, multiplication by $y$, scaling branch, and the two
corrections above).  Until that certificate exists, the general-finite path is
conditional in the same sense as Theorem~\ref{thm:exp-stage}.

\begin{remark}
MPFR probes over the integer, general-finite, and special-value corpora found at
most one ULP error.  That measurement is engineering evidence only and is not
part of the certificates above.
\end{remark}

\section{Source Traceability and Reproducibility}
\label{sec:traceability}

Proof scripts parse constants from the production headers.  A certificate for
duplicated constants would certify the copy, not the implementation.

\begin{enumerate}
  \item \texttt{pow\_double\_proof.py} parses the source arrays from
    \texttt{common\_math\_constants.hpp} and
    \texttt{pow\_impl.hpp}.
  \item \texttt{pow\_float\_proof.py} parses the corresponding arrays from
    \texttt{common\_math\_constants.hpp} and
    \texttt{powf\_impl.hpp}.
  \item Both scripts accept an optional \texttt{--report} flag, which writes a
    JSON record containing tool versions, command lines, source hashes, proof
    inputs, and the certified bounds that were obtained.
  \item For the 80-bit kernel,
    \texttt{tools/proofs/math/power/generate\_powl\_ld80\_tables.py} emits
    \texttt{powl\_ld80\_tables.hpp} from the refit minimax polynomials and the
    binary64 reduction tables, and
    \texttt{validate\_powl\_ld80\_kernel.sollya} certifies the two sup-norm
    bounds of \S\ref{sec:powl-poly} against those same coefficients.  A
    certificate for duplicated constants is a certificate for the duplicate, so
    the two files must agree by construction; the generator and the validator
    carry one and the same coefficient list.
\end{enumerate}

From the repository root, the relevant commands are:

\begin{lstlisting}[language=bash]
python3 tests/proofs/math/power/pow_double_proof.py \
    --report /tmp/pow-double-proof-report.json

python3 tests/proofs/math/power/pow_float_proof.py \
    --report /tmp/pow-float-proof-report.json
\end{lstlisting}

Gappa certificates:

\begin{lstlisting}[language=bash]
gappa tests/proofs/math/power/pow_double_log_eval.g
gappa tests/proofs/math/power/pow_double_exp_eval.g
gappa tests/proofs/math/power/pow_float_log_eval.g
gappa tests/proofs/math/power/pow_float_exp_eval.g
\end{lstlisting}

The 80-bit minimax certificates of \S\ref{sec:powl-poly} are reproduced with:

\begin{lstlisting}[language=bash]
sollya tools/proofs/math/power/validate_powl_ld80_kernel.sollya
\end{lstlisting}

\begin{remark}
Empirical probes may still be useful as engineering diagnostics, but they are
not part of this proof.  The present note is analytic from beginning to end.
\end{remark}

\begin{thebibliography}{9}

\bibitem{ieee754}
IEEE Computer Society,
\textit{IEEE Standard for Floating-Point Arithmetic},
IEEE Std~754-2019, Jul.~2019.
\url{https://doi.org/10.1109/IEEESTD.2019.8766229}

\bibitem{c17}
ISO/IEC,
\textit{Information technology: Programming languages: C},
ISO/IEC~9899:2018 (C17), Jun.~2018.

\bibitem{handbook}
J.-M. Muller, N. Brunie, F. de~Dinechin, C.-P. Jeannerod, M. Joldes,
V. Lef\`evre, G. Melquiond, N. Revol, and S. Torres,
\textit{Handbook of Floating-Point Arithmetic}, 2nd~ed.
Cham: Birkh\"auser/Springer, 2018.
\url{https://doi.org/10.1007/978-3-319-76526-6}

\bibitem{gappa}
M. Daumas and G. Melquiond,
``Certification of bounds on expressions involving rounded operators,''
\textit{ACM Transactions on Mathematical Software},
vol.~37, no.~1, pp.~2:1--2:20, Jan.~2010.
\url{https://doi.org/10.1145/1644001.1644003}

\bibitem{sollya}
S. Chevillard, M. Jolde\c{s}, and C. Lauter,
``Sollya: An environment for the development of numerical codes,''
in \textit{Mathematical Software: ICMS~2010}, ser.~LNCS, vol.~6327.
Berlin: Springer, 2010, pp.~28--31.
\url{https://doi.org/10.1007/978-3-642-15582-6_5}

\bibitem{ziv}
A. Ziv,
``Fast evaluation and rounding of elementary functions,''
\textit{IEEE Transactions on Computers},
vol.~40, no.~7, pp.~843--851, Jul.~1991.
\url{https://doi.org/10.1109/12.83617}

\bibitem{powr}
C. Q. Lauter and V. Lef\`evre,
``An efficient rounding boundary test for \texttt{pow}(x,y) in double
precision,''
\textit{IEEE Transactions on Computers},
vol.~58, no.~2, pp.~197--207, Feb.~2009.
\url{https://doi.org/10.1109/TC.2008.202}
\newline
\url{https://www.vinc17.net/research/papers/ieeetc2009-powr.pdf}

\bibitem{coremath}
A. Sibidanov, P. Zimmermann, and S. Glondu,
``The CORE-MATH Project,''
in \textit{Proceedings of the 29th IEEE Symposium on Computer Arithmetic},
2022.
\url{https://doi.org/10.1109/ARITH54963.2022.00014}

\end{thebibliography}

\end{document}
